%% ============================================================
%% 分片 02：原文第 2 节「${\EUN}=4$ Super Yang-Mills Theory And $S$-Duality」（原文 glshort.tex 起于第 642 行）
%% 由 tools/harvmac2tex.py 从 arXiv 源码（hep-th/0604151v3）自动生成。
%% 只译英文散文；公式、编号、\label、\cite 一律不改。
%% ============================================================

\section{${\EUN}=4$ Super Yang-Mills Theory And $S$-Duality}

\label{sec:review}

 

 In this section, we recall a few properties of
${\EUN}=4$ super Yang-Mills theory and its $S$-duality.

\subsection{Review Of ${\EUN}=4$ Super Yang-Mills Theory}
\label{sec:reviewsym}

As in the original work \cite{ref42}, ${\EUN}=4$ super Yang-Mills is
most easily constructed by dimensional reduction from ten
dimensions. Ten dimensions is the maximum possible dimension for
supersymmetric Yang-Mills theory  by virtue of Nahm's theorem
\cite{ref43}, and for given gauge group $G$, there is a very simple
supersymmetric Lagrangian which moreover is  unique up to the
choice of a few coupling parameters if we ask for a Lagrangian
quadratic in the curvature. In this paper, we always assume $G$ to
be compact and denote its complexification as $G_{\Bbb{C}}$.  This
differs from most other expositions of the geometric Langlands
program.

  
  
  

  

\par\medskip{\it Spacetime Conventions}

 

We begin by describing some conventions. We will work with Lorentz
signature $-++\dots +$ or Euclidean signature $+++\dots +$.
Basically, when  emphasizing questions of topological field
theory, we will use Euclidean signature, but when we want to
stress that the constructions are natural in physically sensible,
unitary quantum field theory, we use Lorentz signature.

We write the metric of ten-dimensional Minkowski space
$\Bbb{R}^{1,9}$ or Euclidean space $\Bbb{R}^{10}$ as
$ds^2=\sum_{I,J=0}^9g_{IJ}\,dx^I\,dx^J=\mp (dx^0)^2+(dx^1)^2+\dots
+(dx^9)^2$. These have symmetry groups $SO(1,9)$ or $SO(10)$,
whose spin representations ${\EUS}^+$ and ${\EUS}^-$ are of rank
16. They are real (and dual to one another) in Lorentz signature,
while in Euclidean signature they are complex conjugate
representations.

The gamma matrices $\Gamma_I,\,I=0,\dots,9$ (which in Lorentz
signature are real) reverse the chirality, mapping ${\EUS}^\pm$ to
${\EUS}^\mp$, and obey the Clifford algebra
$\{\Gamma_I,\Gamma_J\}=2g_{IJ}$. Moreover, the operator
% 式 (2.1)
\setcounter{equation}{0}
\begin{equation}\label{eq:thop}
\bar\Gamma=\Gamma_0\Gamma_1\cdots\Gamma_9
\end{equation}
 acts on
${\EUS}^+$ or ${\EUS}^-$ as multiplication by 1 or $-1$. Because
${\EUS}^+$ and ${\EUS}^-$ are dual, we can regard the $\Gamma$'s
not as maps from ${\EUS}^\pm$ to ${\EUS}^{\mp}$ but as bilinear
maps $\Gamma_I:{\EUS}^+\otimes {\EUS}^+\to \Bbb{R}$ or
$\Gamma_I:{\EUS}^-\otimes{\EUS}^-\to \Bbb{R}$. (In this paragraph,
we assume Lorentz signature; in Euclidean signature, all such maps
are to $\Bbb{C}$, since the spinors are complex.) If
$\beta,\gamma\in {\EUS}^+$, it is conventional to write $\bar\beta
\Gamma_I\gamma$ for the bilinear map $\Gamma_I(\beta,\gamma)$,
which can also be written in components as
$\sum_{a,b=1}^{16}\Gamma_{I\,ab}\beta^a\gamma^b$. (The bar in
$\bar \beta$ is conventional for spinors and should perhaps be
read as transpose, not complex conjugation.) A standard convention
is to define $\Gamma_{I_1I_2\dots I_k}$ to be zero if the indices
$I_1,I_2,\dots,I_k$ are not pairwise distinct and otherwise to
equal the product $\Gamma_{I_1}\Gamma_{I_2}\cdots \Gamma_{I_k}$.
So $\Gamma_{I_1I_2\dots I_k}$ reverses the chirality if $k$ is
odd, and again can be regarded as a bilinear map ${\EUS}^+\otimes
{\EUS}^+\to \Bbb{R}$ (or ${\EUS}^-\otimes {\EUS}^-\to \Bbb{R}$).
These maps are symmetric for $k=1,5,9$ and antisymmetric for
$k=3,7$.  For $\beta,\gamma\in {\EUS}^+$, we again write
$\bar\beta\Gamma_{I_1\dots I_k}\gamma$ for these bilinear maps
$\Gamma_{I_1\dots I_k}(\beta,\gamma)$. For $k$ even, we have
$\Gamma_{I_1I_2\dots I_k}:{\EUS}^\pm\to{\EUS}^\pm$, or
equivalently we have  bilinear maps $\Gamma_{I_1I_2\dots
I_k}:{\EUS}^+\otimes {\EUS}^-\to \Bbb{R}$, again denoted
$\bar\beta\Gamma_{I_1\dots I_k}\gamma$.

\par\medskip{\it Fields, Transformation Laws, And Lagrangian}

The fields of ten-dimensional super Yang-Mills theory are the
gauge field $A$, which is a connection on a $G$-bundle $E$, and a
fermion field $\lambda$ that is a section of ${\EUS}^+\otimes {\rm
ad}(E)$; in other words, $\lambda$ is a positive chirality spinor
field with values in the adjoint representation of $G$. In Lorentz
signature, $\lambda$ is real, since the bundles ${\EUS}^+$ and
${\EUS}^-$ are real.  In Euclidean signature, $\lambda$ is not
real but its complex conjugate does not enter the formalism. The
covariant derivative is $D=d+A$ and the curvature of $A$ is
$F=D^2=dA+A\wedge A$.

We consider $A$ and $\lambda$ to take values in the real Lie
algebra of $G$, which has real structure constants.  This means
that, in a unitary representation of $G$, $A$ and $\lambda$ take
values in antihermitian matrices. This is opposite to the usual
physics convention, but is in accord with the math literature.
(Taking the fields to be antihermitian may look unnatural for
$G=U(1)$, which is the reason for the usual physics convention,
but it avoids unnatural factors of $i$ for general $G$.)  In a
standard set of physics conventions (see p. 4 of \cite{ref44}), the covariant derivative is $D=d-iA'$ with a
hermitian gauge field $A'$.  The relation between our
antihermitian $A$ and this hermitian gauge field $A'$ is thus
% 式 (2.2)
\setcounter{equation}{1}
\begin{equation}\label{eq:boloko}
A=-iA'.
\end{equation}
  The curvature of $A'$ is defined as
$F'=i(D)^2=dA'-i A'\wedge A'$, so % 式 (2.3)
\setcounter{equation}{2}
\begin{equation}\label{eq:obolko}
F=-iF'.
\end{equation}

We use the symbol ``$\Tr$'' to denote an invariant and negative
definite quadratic form on the Lie algebra of $G$.  As we assume
$G$ to be simple, such a quadratic form is unique up to a constant
multiple; we  normalize it so that for $M=S^4$, the characteristic
class ${1\over 8\pi^2}\int_M\Tr\, F\wedge F$ takes arbitrary
integer values. For $G=SU(N)$, this quadratic form can be obtained
as a trace in the $N$-dimensional representation, which motivates
denoting it as $\Tr$. Since our Lie algebras are generated by
antihermitian matrices, $\Tr$ is {\it negative}-definite.

With this understood, the supersymmetry transformation laws and
the Lagrangian of ten-dimensional super Yang-Mills theory can be
described as follows.  The generator of supersymmetry is a
constant spinor $\epsilon$ that takes values in ${\EUS}^+$, and
hence obeys % 式 (2.4)
\setcounter{equation}{3}
\begin{equation}\label{eq:goggo}
\bar\Gamma\epsilon=\epsilon.
\end{equation}
 ($\bar
\Gamma$ was defined in \eqref{eq:thop}; we take $\epsilon$ to be bosonic.)
The supersymmetry transformation generated by $\epsilon$ is
% 式 (2.5)
\setcounter{equation}{4}
\begin{equation}\label{eq:hygo}
\begin{aligned}\delta_S A_I & = i\bar\epsilon\Gamma_I\lambda
\\
                   \delta_S\lambda & = {1\over 2}
                   \Gamma^{IJ}F_{IJ}\epsilon.\\\end{aligned}
\end{equation}

For any field $\Phi$, the symbol $\delta_S \Phi$ is short
for\footnote{In a $\Bbb{Z}_2$-graded algebra, the symbol $[A,B\}$
denotes $AB-(-1)^{AB}BA$.} $\sum_{a=1}^{16}[\epsilon^aQ_a,\Phi\}$,
where $Q_a$, taking values in ${\EUS}^-$, are the sixteen
supersymmetries. The symbol $\delta_S$ stands for supersymmetric
variation. The invariant action is\footnote{This is written in Lorentz
signature, with the usual convention in which the kinetic energy
is positive. For Euclidean signature, one must change the overall
sign of the Lagrangian to make the bosonic part of the action
positive definite.} % 式 (2.6)
\setcounter{equation}{5}
\begin{equation}\label{eq:ygo}
{{\cmmib I}}_{10}={1\over e^2}\int
d^{10}x\,\Tr\left({1\over
2}F_{IJ}F^{IJ}-i\bar\lambda\Gamma^ID_I\lambda\right),
\end{equation}
 with an
arbitrary constant $e$, the gauge coupling.  The verification of
supersymmetry is described in \cite{ref42}. Finally, the conserved
supercurrent that generates the supersymmetries \eqref{eq:hygo}\ is
% 式 (2.7)
\setcounter{equation}{6}
\begin{equation}\label{eq:bygo}
J^I=\half\Tr\, \Gamma^{JK}F_{JK}\Gamma^I\lambda.
\end{equation}

The bosonic symmetry of this theory is not just $SO(1,9)$ but the
``Poincar\'e group'' $\EUP$, which is an extension of $SO(1,9)$ by
the ``translation'' group of $\Bbb{R}^{1,9}$.  This translation
group is isomorphic to $\Bbb{R}^{1,9}$ itself (regarded as an
abelian group), and ${\EUP}$ is an extension % 式 (2.8)
\setcounter{equation}{7}
\begin{equation}\label{eq:ugy}
0\to
\Bbb{R}^{1,9}\to \EUP\to SO(1,9)\to 1.
\end{equation}
  The generators of
$\Bbb{R}^{1,9}$ are called the ``momentum operators'' $P_I$.  The
algebra obeyed by the conserved supercharges $Q_a$, $a=1,\dots,16$
is % 式 (2.9)
\setcounter{equation}{8}
\begin{equation}\label{eq:realify}
\{Q_a,Q_b\}=\sum_{I=1}^{10}\Gamma^I_{ab}P_I.
\end{equation}
 In
addition, the $Q_a$ commute with $P_I$ and transform under
$SO(1,9)$ as ${\EUS}^-$.

\par\medskip{\it Dimensional Reduction To Four Dimensions}

To reduce to four dimensions, we simply take all fields to be
independent of the coordinates $x^4,\dots,x^9$.  The components
$A_I$, $I=0,\dots,3$ describe the four-dimensional gauge field
$A=\sum_{\mu=0}^3A_\mu dx^\mu$, while the components $A_I$, $I\geq
4$, become four-dimensional scalar fields $\phi_i=A_{i+4}$,
$i=0,\dots,5$. The ten-dimensional curvature $F_{IJ}$ has three
types of contribution; depending on whether the number of indices
$I,J$ in the range $4,\dots, 9$ is zero, one, or two, we get a
four-dimensional curvature $F_{\mu\nu}$, a derivative
$D_\mu\phi_i$ of a scalar field, or a commutator $[\phi_i,\phi_j]$
of scalar fields.

The bosonic part of the action, in four dimensions, has all three
types of contribution and becomes % 式 (2.10)
\setcounter{equation}{9}
\begin{equation}\label{eq:loopy}
{{\cmmib
I}}_4={1\over e^2}\int d^4x\,\Tr\,\left({1\over
2}\sum_{\mu,\nu=0}^3F_{\mu\nu}F^{\mu\nu}+\sum_{\mu=0}^3\sum_{i=1}^6D_\mu\phi_i
D^\mu\phi_i +{1\over 2}\sum_{i,j=1}^6[\phi_i,\phi_j]^2\right).
\end{equation}

Together with the part of the action involving fermions, which can
be similarly written in four-dimensional terms, though we will not
do so, this is the essentially unique four-dimensional gauge
theory with the maximal possible supersymmetry.  If $G$ is simple
and if we want a Lagrangian quadratic in derivatives (the case
that leads to a sensible quantum theory), the action is unique
except for the choice of parameter $e$ and, in four dimensions,
another possible parameter that measures the topology of the
$G$-bundle $E$: % 式 (2.11)
\setcounter{equation}{10}
\begin{equation}\label{eq:oopy}
{{\cmmib I}}_{\theta}=-{\theta\over
8\pi^2}\int \Tr \,F\wedge F.
\end{equation}
 This last term is $\theta$ times the
second Chern class or instanton number of the bundle.  The
parameters $e$ and $\theta$ combine into a complex parameter
% 式 (2.12)
\setcounter{equation}{11}
\begin{equation}\label{eq:ingo}
\tau={\theta\over 2\pi}+{4\pi i\over e^2}.
\end{equation}
 As long as
we are on $\Bbb{R}^4$ or any four-manifold $M$ with
$H^2(M,\Bbb{Z})=0$ (for the generalization, in which $\pi_1(G)$
comes into play, see \cite{ref45,ref26}), there is an
elementary symmetry $\tau\to\tau+1$, which expresses the fact that
$(1/8\pi^2)\int\Tr F\wedge F$ is integer-valued, and that in
quantum mechanics one only cares about the action modulo an
integer multiple of $2\pi$. Equivalently, $\theta$ is an angular
variable, with $\theta\cong\theta+2\pi$.

The $SO(1,9)$ (or $SO(10)$) symmetry in ten dimensions becomes,
after dimensional reduction to four dimensions, $SO(1,3)\times
SO(6)$ (or $SO(4)\times SO(6)$). Allowing for the presence of
spinors in the theory, the symmetry is really $Spin(1,9)$ reduced
to $Spin(1,3)\times Spin(6)$ (or $Spin(10)$ reduced to
$Spin(4)\times Spin(6)$).  The group $Spin(6)$ is isomorphic to
$SU(4)$ and is known as the ``$R$ symmetry group'' of the theory.
We will call it $SU(4)_{{\cal R}}$.

\def\btwo{{\bf 2}}
\def\bone{{\bf 1}}
\def\b4{{\bf 4}}
\def\bb4{\bar{\bf 4}}
 The chirality condition $\bar\Gamma\epsilon=\epsilon$ in ten
dimensions becomes in four dimensions % 式 (2.13)
\setcounter{equation}{12}
\begin{equation}\label{eq:yug}
\hat\Gamma
\Gamma'\epsilon=\epsilon,
\end{equation}
 where
$\hat\Gamma=\Gamma_0\Gamma_1\Gamma_2\Gamma_3$ measures the
$Spin(1,3)$ chirality and $\Gamma'=\Gamma_4\Gamma_5\dots\Gamma_9$
measures the $Spin(6)$ chirality. $\hat \Gamma$ and $\Gamma'$ have
eigenvalues $\pm i$; \eqref{eq:yug}\ means that the eigenvalue of $\Gamma'$
is minus that of $\hat \Gamma$. The two eigenvalues of
$\hat\Gamma$ distinguish the two spin representations  of
$Spin(1,3)$, while the eigenvalues of $\Gamma'$ similarly label
the spin representations of $Spin(6)$. The complexification of
$Spin(1,3)$ is $SL(2)\times SL(2)$ and the two spin
representations correspond to the representations $({\bf 2},{\bf
1})$ and $({\bf 1},{\bf 2})$ of $SL(2)\times SL(2)$ (here $({\bf
2},{\bf 1})$ is the two-dimensional representation of the first
$SL(2)$ tensored with the trivial one-dimensional representation
of the second $SL(2)$, and vice-versa for $(\bone,\btwo)$). The
spin representations of $Spin(6)$ are the defining
four-dimensional representation of $SU(4)_{{\cal R}}$ and its
dual; we denote them as ${\bf 4}$ and $\bar{\bf 4}$. We pick
orientations so that $\hat \Gamma$ acts as $i^{-1}$ or $i$ on the
$({\bf 2},{\bf 1})$ and $({\bf 1},{\bf 2})$, respectively, and
$\Gamma'$ acts as $i$ and $i^{-1}$ on the $\bar {\bf 4}$ and ${\bf
4}$.
 So
\eqref{eq:yug}\ means that the four-dimensional supersymmetries transform
under $Spin(1,3)\times Spin(6)\sim SL(2)\times SL(2)\times
Spin(6)$ as % 式 (2.14)
\setcounter{equation}{13}
\begin{equation}\label{eq:jugo}
(\btwo,\bone,\bb4)\oplus (\bone,\btwo,\b4).
\end{equation}

The fermion fields $\lambda$ transform the same way under
$Spin(1,3)\times Spin(6)$.

If we write $\bar Q_{AX}$, $A=1,2$, $X=1,\dots, 4$ for the
supersymmetries of type $({\bf 2},\bone,\bar{\bf 4})$, and
similarly $ Q_{\dot A }^Y$, $\dot A=1,2$, $Y=1,\dots,4$ for those
of type $(\bone,{\bf 2},{\bf  4})$, then the algebra generated by
the supersymmetries is the reduction of \eqref{eq:realify},
% 式 (2.15)
\setcounter{equation}{14}
\begin{equation}\label{eq:ugo}
\begin{aligned}\{\bar Q_{AX}, Q_{\dot A}^
Y\}&=\delta_X{}^Y\sum_{\mu=0}^3\Gamma^\mu_{A\dot A}P_\mu\\
\{Q,Q\}=\{\bar Q,\bar Q\}&=0,\\\end{aligned}
\end{equation}
 where now the four momentum
operators $P_\mu$ generate the translations of $\Bbb{R}^{1,3}$.
With suitable boundary conditions, additional terms appear  \cite{ref13}\
on the right hand side of \eqref{eq:ugo}\ (they are related to the extra six
momenta that were dropped in going to four dimensions, and their
magnetic duals).  They will make a brief appearance at the end of
this section.
